\documentclass[12pt,a4paper]{article}
\usepackage[utf8]{inputenc}
\usepackage[spanish]{babel}
\usepackage{graphicx}
\usepackage{listings}
\usepackage{xcolor}
\usepackage[margin=2.5cm]{geometry}
\usepackage{hyperref}

\definecolor{codebg}{RGB}{245,245,245}
\definecolor{codegreen}{RGB}{0,128,0}
\definecolor{codepurple}{RGB}{128,0,128}
\definecolor{codeblue}{RGB}{0,0,200}

\lstset{
    language=Java,
    backgroundcolor=\color{codebg},
    basicstyle=\ttfamily\footnotesize,
    keywordstyle=\color{codeblue}\bfseries,
    stringstyle=\color{codegreen},
    commentstyle=\color{codegreen}\itshape,
    numberstyle=\tiny\color{gray},
    numbers=left,
    numbersep=5pt,
    breaklines=true,
    frame=single,
    rulecolor=\color{gray},
    tabsize=4,
    showstringspaces=false,
    literate={á}{{\'a}}1 {é}{{\'e}}1 {í}{{\'i}}1 {ó}{{\'o}}1 {ú}{{\'u}}1 {ñ}{{\~n}}1
}

\title{\textbf{Antidoping -- Iteración 1} \\ \large Ingeniería del Software 2025--2026}
\author{Ibai Martinez}
\date{\today}

\begin{document}

\maketitle

Para cada CU nuevo de la iteración 1 se indica: la clase de presentación, los métodos de lógica de negocio añadidos (con la clase donde se implementan) y las clases del dominio utilizadas.

% ============================================================
\section{CU: Registrar Usuario}
% ============================================================

\textbf{Clase de presentación:} \texttt{gui.RegisterGUI}

Formulario con email, nombre, contraseña y selector de rol (Comprador/Vendedor). Según el rol muestra campo de dirección de envío o cuenta bancaria.

\subsection{Métodos de lógica de negocio}

Implementados en \texttt{businessLogic.BLFacadeImplementation}:

\begin{lstlisting}
// Registrar comprador
public Buyer registerBuyer(String email, String name,
    String password, String shippingAddress)
    throws UserAlreadyExistsException, InvalidEmailException,
           InvalidFieldException;

// Registrar vendedor
public Seller registerSeller(String email, String name,
    String password, String bankAccount)
    throws UserAlreadyExistsException, InvalidEmailException,
           InvalidFieldException;
\end{lstlisting}

Ambos delegan en \texttt{dataAccess.DataAccess} (\texttt{createBuyer()} y \texttt{createSeller()} respectivamente), que validan campos, verifican email único y persisten la entidad.

\subsection{Clases del dominio}

\begin{itemize}
    \item \texttt{domain.User} -- Clase base (email, nombre, contraseña).
    \item \texttt{domain.Buyer} -- Hereda de User. Añade \texttt{shippingAddress}.
    \item \texttt{domain.Seller} -- Hereda de User. Añade \texttt{bankAccount}.
\end{itemize}

% ============================================================
\section{CU: Iniciar Sesión}
% ============================================================

\textbf{Clase de presentación:} \texttt{gui.LoginGUI}

Pantalla inicial de la aplicación. Campos de email y contraseña. Según el tipo de usuario autenticado abre \texttt{BuyerMainGUI} o \texttt{MainGUI}.

\subsection{Métodos de lógica de negocio}

Implementado en \texttt{businessLogic.BLFacadeImplementation}:

\begin{lstlisting}
// Autenticar usuario
public User login(String email, String password);
\end{lstlisting}

Delega en \texttt{DataAccess.login()}, que busca al usuario por email y verifica la contraseña. Retorna el \texttt{User} (puede ser \texttt{Buyer} o \texttt{Seller}) o \texttt{null}.

\subsection{Clases del dominio}

\begin{itemize}
    \item \texttt{domain.User} -- Clase base buscada en BD por email.
    \item \texttt{domain.Buyer} -- Subtipo, se detecta con \texttt{instanceof}.
    \item \texttt{domain.Seller} -- Subtipo, se detecta con \texttt{instanceof}.
\end{itemize}

% ============================================================
\section{CU: Aceptar Oferta}
% ============================================================

\textbf{Clase de presentación:} \texttt{gui.AcceptOfferGUI}

Tabla con ofertas disponibles (número, título, descripción, precio, vendedor). Campo opcional para precio negociado. Botones de aceptar y refrescar.

\subsection{Métodos de lógica de negocio}

Implementados en \texttt{businessLogic.BLFacadeImplementation}:

\begin{lstlisting}
// Aceptar una oferta con posible negociación
public AcceptedOffer acceptOffer(String buyerEmail,
    Integer saleNumber, Float negotiatedPrice)
    throws InvalidPriceException;

// Obtener ofertas disponibles
public List<Sale> getAvailableSalesForBuyer(Date pubDate);
\end{lstlisting}

\texttt{acceptOffer()} delega en \texttt{dataAccess.DataAccess.acceptOffer()}, que busca Buyer y Sale, valida el precio negociado (positivo y $\leq$ precio original), crea un \texttt{AcceptedOffer} y lo persiste.

\texttt{getAvailableSalesForBuyer()} delega en \texttt{dataAccess.DataAccess}, que ejecuta \texttt{SELECT s FROM Sale s}.

\subsection{Clases del dominio}

\begin{itemize}
    \item \texttt{domain.Buyer} -- Comprador que acepta la oferta.
    \item \texttt{domain.Sale} -- Oferta aceptada. Su precio se usa para validar la negociación.
    \item \texttt{domain.AcceptedOffer} -- Registro de aceptación (buyer, sale, precio negociado, fecha).
    \item \texttt{domain.Seller} -- Accedido vía \texttt{Sale.getSeller()} para mostrar el vendedor.
\end{itemize}

% ============================================================
\section{CU: Visualizar Ofertas Aceptadas}
% ============================================================

\textbf{Clase de presentación:} \texttt{gui.ViewAcceptedOffersGUI}

Tabla con las ofertas aceptadas por el comprador: título, precio, comprador, email y fecha de aceptación. Botón de refrescar.

\subsection{Métodos de lógica de negocio}

Implementado en \texttt{businessLogic.BLFacadeImplementation}:

\begin{lstlisting}
// Obtener ofertas aceptadas por un comprador
public List<AcceptedOffer> getAcceptedOffersByBuyer(
    String buyerEmail);
\end{lstlisting}

Delega en \texttt{dataAccess.DataAccess.getAcceptedOffersByBuyer()}, que ejecuta:
\texttt{SELECT a FROM AcceptedOffer a WHERE a.buyer.email = :email}.

\subsection{Clases del dominio}

\begin{itemize}
    \item \texttt{domain.AcceptedOffer} -- Entidad consultada (fecha, precio, referencias).
    \item \texttt{domain.Buyer} -- Accedido vía \texttt{AcceptedOffer.getBuyer()}.
    \item \texttt{domain.Sale} -- Accedido vía \texttt{AcceptedOffer.getSale()}.
\end{itemize}

\end{document}
